% GENERATED FILE. Edit canonical lessons, then run npm run generate:books.
% canonical-source-hash: fnv1a64:260ee7afd7dbac46
% canonical-lessons: IT-C328-impero, IT-C328-imperatore, IT-C328-sovrano, IT-C328-governare, IT-C328-fondare, IT-R328-ripasso-321-328

\chapter{Imperi e sovrani}
\label{ch:it-imperi-e-sovrani}
\hlchaptermodality{328}

\begin{chapteropening}
\textbf{By the end of this chapter:} \emph{I can talk about empires: the emperor, the ruler, governing and founding a city.}

Complete the last lesson of chapter 328: I can talk about empires: the emperor, the ruler, governing and founding a city.
\end{chapteropening}

\section[l'impero]{l'impero — the empire}

\label{lesson:IT-C328-impero}



\begin{quote}
Before the new one: say the Italian for the shield, then the Italian for to conquer.
\end{quote}

\subsection*{You'll want to know: l'impero}
\textbf{l'impero} — \textquotedblleft{}the empire\textquotedblright{}.

\begin{quote}
L'impero è durato cinque secoli.
\end{quote}

\begin{quote}
\emph{The empire lasted five centuries.}
\end{quote}

Built on Latin \emph{imperium}, ‘command’.

\begin{cousinweb}
\emph{impero} is inherited from Latin \emph{imperium} \textquotedblleft{}command, power, empire\textquotedblright{} (EtymDB), from \emph{imperāre} \textquotedblleft{}to command\textquotedblright{}.
\end{cousinweb}

\begin{grammarlens}[title={What you've built}]
One more word for telling, explaining and describing at B1.
\end{grammarlens}

\subsection*{Your turn}
\begin{itemize}
\raggedright
  \item \emph{Say it:} \emph{l'impero}
  \item \emph{Say it:} \emph{l'impero}, once more
  \item \emph{Say it:} a sentence of your own with \emph{l'impero}
  \item \emph{Recall:} say the Italian for the shield, then the Italian for to conquer, then say \emph{l'impero} again
\end{itemize}

\begin{tcolorbox}[breakable,colback=teal!4,colframe=teal!35!black,title={Before you move on}]
What does l'impero mean? (\textquotedblleft{}The empire\textquotedblright{}.) Say it once more.
\end{tcolorbox}
\section[l'imperatore]{l'imperatore — the emperor}

\label{lesson:IT-C328-imperatore}



\begin{quote}
Before the new one: say the Italian for to conquer, then the Italian for the empire.
\end{quote}

\subsection*{You'll want to know: l'imperatore}
\textbf{l'imperatore} — \textquotedblleft{}the emperor\textquotedblright{}.

\begin{quote}
L'imperatore ha costruito un grande palazzo.
\end{quote}

\begin{quote}
\emph{The emperor built a great palace.}
\end{quote}

An empress is \emph{l'imperatrice}.

\begin{cousinweb}
\emph{imperatore} is inherited from Latin \emph{imperātor} \textquotedblleft{}a commander\textquotedblright{} (EtymDB).
\end{cousinweb}

\begin{grammarlens}[title={What you've built}]
One more word for telling, explaining and describing at B1.
\end{grammarlens}

\subsection*{Your turn}
\begin{itemize}
\raggedright
  \item \emph{Say it:} \emph{l'imperatore}
  \item \emph{Say it:} \emph{l'imperatore}, once more
  \item \emph{Say it:} a sentence of your own with \emph{l'imperatore}
  \item \emph{Recall:} say the Italian for to conquer, then the Italian for the empire, then say \emph{l'imperatore} again
\end{itemize}

\begin{tcolorbox}[breakable,colback=teal!4,colframe=teal!35!black,title={Before you move on}]
What does l'imperatore mean? (\textquotedblleft{}The emperor\textquotedblright{}.) Say it once more.
\end{tcolorbox}
\section[il sovrano]{il sovrano — the sovereign, the ruler}

\label{lesson:IT-C328-sovrano}



\begin{quote}
Before the new one: say the Italian for the empire, then the Italian for the emperor.
\end{quote}

\subsection*{You'll want to know: il sovrano}
\textbf{il sovrano} — \textquotedblleft{}the sovereign, the ruler\textquotedblright{}.

\begin{quote}
Il sovrano ha lasciato il palazzo.
\end{quote}

\begin{quote}
\emph{The ruler left the palace.}
\end{quote}

Also an adjective: \emph{uno stato sovrano}.

\begin{cousinweb}
\emph{sovrano} is inherited from Vulgar Latin \emph{superānus} (EtymDB), built on \emph{super} \textquotedblleft{}above\textquotedblright{}.
\end{cousinweb}

\begin{grammarlens}[title={What you've built}]
One more word for telling, explaining and describing at B1.
\end{grammarlens}

\subsection*{Your turn}
\begin{itemize}
\raggedright
  \item \emph{Say it:} \emph{il sovrano}
  \item \emph{Say it:} \emph{il sovrano}, once more
  \item \emph{Say it:} a sentence of your own with \emph{il sovrano}
  \item \emph{Recall:} say the Italian for the empire, then the Italian for the emperor, then say \emph{il sovrano} again
\end{itemize}

\begin{tcolorbox}[breakable,colback=teal!4,colframe=teal!35!black,title={Before you move on}]
What does il sovrano mean? (\textquotedblleft{}The sovereign, the ruler\textquotedblright{}.) Say it once more.
\end{tcolorbox}
\section[governare]{governare — to govern, to rule}

\label{lesson:IT-C328-governare}



\begin{quote}
Before the new one: say the Italian for the emperor, then the Italian for the sovereign.
\end{quote}

\subsection*{You'll want to know: governare}
\textbf{governare} — \textquotedblleft{}to govern, to rule\textquotedblright{}.

\begin{quote}
Ha governato il paese per dieci anni.
\end{quote}

\begin{quote}
\emph{He governed the country for ten years.}
\end{quote}

\emph{Il governo} — the government.

\begin{cousinweb}
\emph{governare} is inherited from Latin \emph{gubernāre} \textquotedblleft{}to steer a ship\textquotedblright{} (EtymDB), from Greek \emph{kybernân}.
\end{cousinweb}

\begin{grammarlens}[title={What you've built}]
One more word for telling, explaining and describing at B1.
\end{grammarlens}

\subsection*{Your turn}
\begin{itemize}
\raggedright
  \item \emph{Say it:} \emph{governare}
  \item \emph{Say it:} \emph{governare}, once more
  \item \emph{Say it:} a sentence of your own with \emph{governare}
  \item \emph{Recall:} say the Italian for the emperor, then the Italian for the sovereign, then say \emph{governare} again
\end{itemize}

\begin{tcolorbox}[breakable,colback=teal!4,colframe=teal!35!black,title={Before you move on}]
What does governare mean? (\textquotedblleft{}To govern, to rule\textquotedblright{}.) Say it once more.
\end{tcolorbox}
\section[fondare]{fondare — to found, to establish}

\label{lesson:IT-C328-fondare}



\begin{quote}
Before the new one: say the Italian for the sovereign, then the Italian for to govern.
\end{quote}

\subsection*{You'll want to know: fondare}
\textbf{fondare} — \textquotedblleft{}to found, to establish\textquotedblright{}.

\begin{quote}
Secondo la leggenda, due fratelli hanno fondato la città.
\end{quote}

\begin{quote}
\emph{According to the legend, two brothers founded the city.}
\end{quote}

\emph{Fondato nel 1900} — founded in 1900.

\begin{cousinweb}
\emph{fondare} is inherited from Latin \emph{fundāre} \textquotedblleft{}to lay the bottom, to found\textquotedblright{} (EtymDB), from \emph{fundus} \textquotedblleft{}bottom\textquotedblright{}.
\end{cousinweb}

\begin{grammarlens}[title={What you've built}]
One more word for telling, explaining and describing at B1.
\end{grammarlens}

\subsection*{Your turn}
\begin{itemize}
\raggedright
  \item \emph{Say it:} \emph{fondare}
  \item \emph{Say it:} \emph{fondare}, once more
  \item \emph{Say it:} a sentence of your own with \emph{fondare}
  \item \emph{Recall:} say the Italian for the sovereign, then the Italian for to govern, then say \emph{fondare} again
\end{itemize}

\begin{tcolorbox}[breakable,colback=teal!4,colframe=teal!35!black,title={Before you move on}]
What does fondare mean? (\textquotedblleft{}To found, to establish\textquotedblright{}.) Say it once more.
\end{tcolorbox}
\section[(dialogue)]{Ripasso dei capitoli 321–328}

\label{lesson:IT-R328-ripasso-321-328}



\begin{quote}
Say the Italian for the nativity scene and the panettone.
\end{quote}

\subsection*{Your turn}
Say these aloud:

\begin{itemize}
\raggedright
  \item the nativity scene, the panettone, the big festive dinner, to decorate, the tombola
  \item the fireworks, the firework bang, the bonfire, the lantern, to dance
  \item once upon a time, the story, the fable, to narrate, the narrator
  \item the hero, the legend, the myth, the ending, the moral
  \item the prince, the princess, the queen, the kingdom, to reign
  \item the dragon, the witch, the fairy, the wizard, the treasure
  \item the battle, the knight, the sword, the shield, to conquer
  \item the empire, the emperor, the sovereign, to govern, to found
\end{itemize}

\begin{tcolorbox}[breakable,colback=teal!4,colframe=teal!35!black,title={Before you move on}]
The nativity scene? (\textbf{il presepe}.) The panettone? (\textbf{il panettone}.) The big festive dinner? (\textbf{il cenone}.) To decorate? (\textbf{addobbare}.) The tombola? (\textbf{la tombola}.) The fireworks? (\textbf{i fuochi d'artificio}.) The firework bang? (\textbf{il botto}.) The bonfire? (\textbf{il falò}.) The lantern? (\textbf{la lanterna}.) To dance? (\textbf{danzare}.) Once upon a time? (\textbf{c'era una volta}.) The story? (\textbf{il racconto}.) The fable? (\textbf{la favola}.) To narrate? (\textbf{narrare}.) The narrator? (\textbf{il narratore}.) The hero? (\textbf{l'eroe}.) The legend? (\textbf{la leggenda}.) The myth? (\textbf{il mito}.) The ending? (\textbf{il finale}.) The moral? (\textbf{la morale}.) The prince? (\textbf{il principe}.) The princess? (\textbf{la principessa}.) The queen? (\textbf{la regina}.) The kingdom? (\textbf{il regno}.) To reign? (\textbf{regnare}.) The dragon? (\textbf{il drago}.) The witch? (\textbf{la strega}.) The fairy? (\textbf{la fata}.) The wizard? (\textbf{il mago}.) The treasure? (\textbf{il tesoro}.) The battle? (\textbf{la battaglia}.) The knight? (\textbf{il cavaliere}.) The sword? (\textbf{la spada}.) The shield? (\textbf{lo scudo}.) To conquer? (\textbf{conquistare}.) The empire? (\textbf{l'impero}.) The emperor? (\textbf{l'imperatore}.) The sovereign? (\textbf{il sovrano}.) To govern? (\textbf{governare}.) To found? (\textbf{fondare}.)
\end{tcolorbox}
